% Source printed pp. 506-507 / PDF pp. 514-515.
\subsection{一般型簇的 KSB 模空间}\label{src:10.1.3-subsection}
KSB 模空间 $M_{n,V}^{\mathrm{KSB}}$ 将 $\overline M_g$ 推广到高维一般型簇. 在定理~\hyperref[thm:A]{A} 中, 我们假设亏格 $g$ 至少为 $2$, 这对应于要求曲线为一般型. 在维数 $n>1$ 时, 一般型簇通常未必具有典范极化. 我们要研究的模问题是双有理等价意义下的一般型簇; 由于每个一般型簇都具有典范模型 \cite{BCHM10}, 这个模问题等价于研究同构意义下的典范极化簇. 射影模空间 $M_{n,V}^{\mathrm{KSB}}$ 的存在, 是过去几十年来极小模型纲领 (MMP), 奇点理论和模理论诸多进展的巅峰成果. 这一纲领的开创者包括 Alexeev, Hacon, Koll\'ar, Kov\'acs, McKernan, Mori, Shepherd-Barron, Viehweg 和 Xu, 另有许多数学家作出了重要贡献. Koll\'ar 最近的奠基性工作 \cite{Kol23} 完成了这一纲领. 平易的入门介绍见 \cite{Kov09}.

早期用 GIT 构造一般型簇的模空间的尝试仅取得了部分成功. Gieseker 通过极其艰巨的 GIT 稳定性分析, 验证了多重典范嵌入渐近 Hilbert 稳定, 从而构造了光滑典范极化曲面 $(S,K_S)$ 的拟射影模空间 \cite{Gie77b}. 另一方面, \cite{WX14} 中的例子表明, Chow 半稳定性一般不会在渐近意义下稳定下来.

Koll\'ar 与 Shepherd-Barron 的里程碑论文 \cite{KSB88} 研究了典范极化曲面的模空间, 并作出了关键观察: $n+1$ 维的 MMP 能给出令人满意的 $n$ 维模理论. 基本思路是以固有性的赋值判别法为指引, 确定光滑一般型簇的紧化应允许哪些奇异簇. 这与我们利用曲面的双有理几何证明特征 $0$ 下曲线的稳定约化的方法 (定理~\bookref{6.5.1}) 相仿: 给定穿孔圆盘上的光滑族 $\mathcal X^*\to\Delta^*$, 首先找到延拓 $\mathcal X\to\Delta$, 其中心纤维可以任意奇异; 接着取奇点消解 $\widetilde{\mathcal X}\to\mathcal X$; 最后在 $\widetilde{\mathcal X}$ 上相对于 $\Delta$ 运行 MMP, 得到新族 $\mathcal X'\to\Delta$. 中心纤维 $\mathcal X'_0$ 是具有 $\mathbb Q$-Gorenstein 半对数典范奇点 (slc) 且典范除子丰沛的射影簇.

这促使我们将\newterm{稳定簇}{stable variety}定义为具有 ($\mathbb Q$-Gorenstein) \newterm{半对数典范奇点}{semi-log canonical singularities}且 $\omega_X$ 丰沛的射影簇 $X$. 即使已确定稳定簇是``正确''的奇异簇类别, 构造 KSB 模空间仍有许多额外的细节和困难. 定义稳定簇族 $\mathcal X\to S$ 时, 仅要求 $\mathcal X\to S$ 固有平坦且具有稳定几何纤维是不够的. 例如, 存在曲面族 $\mathcal X\to S$, 每条纤维都是 $\mathbb Q$-Gorenstein, 但 $\omega_{\mathcal X/S}$ 并非 $\mathbb Q$-Cartier, 而且存在点 $s,t\in S$ 满足 $K_{\mathcal X_s}^2\neq K_{\mathcal X_t}^2$; 见 \cite[例 1.42-7]{Kol23}. 这些病态由 \newterm{Koll\'ar 条件}{Koll\'ar's condition}排除: 自反包 $\omega_{\mathcal X/S}^{[m]}:=(\omega_{\mathcal X/S}^{\otimes m})^{\vee\vee}$ 在 $S$ 上平坦且与基变换 $T\to S$ 交换. (这个条件也可用 Abramovich--Hassett 的扭曲簇叠来表述 \cite{AH11}.)

KSB 模叠 $\mathcal M_{n,V}^{\mathrm{KSB}}$ 在概形 $S$ 上的对象是平坦固有有限表现态射 $\mathcal X\to S$, 其几何纤维为维数 $n$, 体积 $V$ 的稳定簇, 且满足 Koll\'ar 条件. 族上的 Koll\'ar 条件还有若干变体, 同样产生性质良好的模空间: 它们在拓扑上与 $M_{n,V}^{\mathrm{KSB}}$ 相同, 但具有不同的概形结构; 见 \cite[\S6]{Kol23}. 叠 $\mathcal M_{n,V}^{\mathrm{KSB}}$ 参数化维数为 $n$, 体积为 $V$ 的稳定簇. 回忆 $n$ 维簇 $X$ 的体积定义为最高次交数 $K_X^n$. 对稳定簇而言, 维数和体积能确定 Hilbert 多项式.

与定理~\bookref{4.1.17} 一样, 可将模叠表为 Hilbert 概形的商, 从而建立 $\mathcal M_{n,V}^{\mathrm{KSB}}$ 的代数性. 这需要证明一个精细的事实: 多重典范嵌入的稳定簇在 Hilbert 概形中是局部闭的; 其证明利用了 Koll\'ar 条件 \cite[推论 24]{Kol09}. 稳定性的有界性, 对光滑簇由 \cite{Mat72} 建立, 对稳定曲面由 \cite{Ale93} 建立, 一般情形则由 \cite{HMX18} 建立. 按上述思路利用 MMP 和赋值判别法可得固有性 \cite[定理 2.51]{Kol23}. 固有粗模空间 $M_{n,V}^{\mathrm{KSB}}$ 的存在由 Keel--Mori 定理 (\bookref{5.4.6}) 得到, 射影性则利用 Koll\'ar 丰沛性判别法证明 \cite{KP17}.

还有\newterm{稳定对}{stable pair} $(X,D)$ 的概念, 由 slc 射影簇 $X$ 与 $\mathbb Q$-除子 $D$ 组成, 满足 $K_X+D$ 丰沛, 最初由 Alexeev 引入 \cite{Ale93}. 定义稳定对的族和构造模空间时还有额外的细节; 详细论述见 \cite{Kol23}. 结论是: 对任意元组 $a\in\mathbb Q^n$, 存在射影 KSBA 模空间 $M_{n,a,V}^{\mathrm{KSBA}}$, 参数化维数为 $n$, 体积为 $V=(K_X+D)^n$ 的稳定对 $(X,D=\sum a_iD_i)$.

KSB 与 KSBA 模空间一般不像 $\overline M_{g,n}$ 那样性质良好. 它们可能极其奇异, 非约化, 并有许多维数不同的不可约分支. 即使对于光滑典范极化曲面, 模空间也满足 Murphy 定律 \cite{Vak06}.
